% !TEX root = final_report.tex
\documentclass[final_report.tex]{subfiles}
\setcounter{chapter}{2}

\begin{document}
\ourchapter{RoPART: design and implementation}\label{ch:ropart}

\section{The \texorpdfstring{\SEtwo}{SE(2)} target}\label{sec:se2-target}
\stub{Extending the per-pair target from \dxy{} to \dxy{} plus \dphi{}, and the
structural invariants it must satisfy: identity on the diagonal, symmetry in
\(\cos\dphi\), antisymmetry in \(\sin\dphi\). Explains the build order --- one
channel group at a time, simplest first --- so any downstream effect is attributable.}

\section{Angle representation and loss}\label{sec:angle-representation}
\stub{Why \((\cos\dphi, \sin\dphi)\) regression rather than a raw angle (the
\(\pm\pi\) wrap) or discrete bins. The loss, and the \(w_\varphi\) group weighting
that later turns out to matter more than expected.}
To represent rotation, a simple direct angle regression may be used. This is a poor neural target: an angle is a point on the unit circle, so any scalar parametrisation is discontinuous at the $\pm\pi$ wrap-around, where two near-identical orientations sit at opposite ends of the range. The network must then approximate a discontinuous function, which yields large gradients at the seam and penalises errors that reflect no real geometric difference; Zhou et al.~\cite{rotcont} show such discontinuous representations are harder to learn. Encoding the angle on the unit circle as $(\cos\Delta\varphi,\sin\Delta\varphi)$ removes the seam and is therefore a better target, supervised either by an MSE on the two components or by a cosine loss $1-\cos(\Delta\hat\varphi-\Delta\varphi)$. This representation also fixes how the intrinsic metrics are computed: the predicted angle is recovered from the network's $(\cos,\sin)$ outputs, the angular error is measured as the difference from the true angle wrapped around the circle into $(-\pi,\pi]$, and the negative-symmetry and compositional-consistency checks are likewise evaluated on the circle rather than on raw scalar differences.

Predicting several relative quantities from a shared backbone is a multi-task regression. This means that channels with different units and convergence rate may interfere and a single fixed loss weighting may prove hard to tune and may not be optimal throughout training. Indeed, in PART~\cite{part} itself adding the relative scale and aspect ratio terms delayed the convergence of translation terms relative to the baseline. One established solution is to make the weighting adaptive. Kendall et al.~\cite{kendall} propose weighting each task by a learned homoscedastic uncertainty which automatically downweights noisier tasks.


\section{Patch position convention}\label{sec:position-convention}
\stub{Centred versus corner anchors, and why the centre is the rotation-invariant
choice that keeps translation decoupled from \dphi{}. Notes the reference-local
frame as the deferred alternative.}

\section{Rotating pixels}\label{sec:rotating-pixels}
\stub{The crop--rotate--centre-crop--resize pipeline, which avoids empty corners.
The two rotation sets: continuous bounded rotation, and the discrete quad set of
exact \texttt{rot90} turns.}

\section{Interpolation controls by construction}\label{sec:interpolation-controls}
\stub{The resampling-blur signature correlated with \(\varphi\), and the control
arms built to break it: raw, supersampled, and matched or randomised blur. Framed
as a false-positive generator that any positive result must be tested against.}

\section{The relative head}\label{sec:relative-head}
\stub{The cross-attention relative encoder, the two query types, and the baseline
that would not learn under the upstream default --- diagnosed with the single-batch
overfit harness.}

\section{Implementation}\label{sec:implementation}
\stub{The \texttt{ropart} package: flag-gated diffs against the upstream path, the
test suite, and device-agnostic execution across MPS and CUDA.}

\section{Experimental infrastructure}\label{sec:infrastructure}
\stub{Cluster execution, atomic checkpointing and supervised restart, experiment
tracking, and what it takes to make a multi-day run survive a dead GPU box.}

\section{Protocol and metrics}\label{sec:protocol-and-metrics}
\stub{Datasets and model configurations. The mean-predictor floor and its
derivation, the four intrinsic gates --- identity, negative symmetry, compositional
consistency, translation --- and angular MAE. Every later result is quoted against
these, so they are defined once, here.}
\paragraph{Intrinsic metrics.} These characterise the orientation channel directly,
computed on the unit circle to respect the $\pm\pi$ wrap-around:
\begin{itemize}
  \item \emph{Angular error} on $\Delta\varphi$, recovered from the predicted
  $(\cos,\sin)$ outputs and wrapped into $(-\pi,\pi]$. The reference null is the
  mean-predictor floor: for approximately uniform $\Delta\varphi$, predicting
  $(0,0)$ gives a $(\cos,\sin)$ MSE of exactly $1.0$, so beating this floor by a
  clear margin is the minimal evidence that anything is being learned.
  \item \emph{Rotational negative symmetry}: predictions should satisfy
  $\Delta\varphi(i\!\to\!j) = -\Delta\varphi(j\!\to\!i)$, the rotational analogue of
  the translational symmetry already confirmed in Phase~0.
  \item \emph{Compositional consistency}: relative orientations should compose
  consistently around triples of patches.
\end{itemize}
A positive intrinsic result is angular error well below the mean-predictor floor with
the symmetry and consistency invariants holding --- \emph{and}, critically, surviving
a matched-interpolation control (the same resampling applied to $\varphi = 0$
patches) rather than reflecting a resampling artefact.

% =====================================================================
%  WORKING SECTION --- delete before submission.
%  Equation bank for Chapter 3. Every equation below is tagged with the
%  section it is intended for; move it into that section's prose and
%  delete it from here. Each block states what the equation asserts and,
%  where it matters, which line of the implementation it was read off.
% =====================================================================
\section{Working notes: equation bank}\label{sec:eqbank}

\stub{Scaffolding, not prose. Notation follows the source paper as summarised in
\cref{sec:part}: $N$ patches of side $P$ from an $S\times S$ image, token width $d$,
patch embeddings $X'$, per-pair target $\theta$ and prediction $\hat\theta$. Ordered
pairs are written $(i,j)$ with $i$ the \emph{reference} and $j$ the \emph{target},
matching PART's $(\mathrm{ref},\mathrm{tgt})$. Every equation here was read off the
code that produced the results, and departures from the paper's presentation are
flagged as such. A few blocks are marked \textbf{NOT IMPLEMENTED}: they describe an
alternative discussed in the text but never run, and must not be written as though
they were.}

Two operators are used throughout. The planar rotation matrix
\begin{equation}\label{eq:rotmat}
  R(\varphi) \;=\;
  \begin{bmatrix} \cos\varphi & -\sin\varphi \\ \sin\varphi & \phantom{-}\cos\varphi \end{bmatrix},
\end{equation}
and the wrap-to-$(-\pi,\pi]$ operator, which is what makes every angular quantity in
this report well defined,
\begin{equation}\label{eq:wrap}
  \operatorname{wrap}(\alpha) \;=\; \bigl((\alpha + \pi) \bmod 2\pi\bigr) - \pi
  \;\in\; (-\pi, \pi].
\end{equation}
In code \cref{eq:wrap} is the line \py{(pred - true + math.pi) \% (2*math.pi) - math.pi},
which appears identically in \py{engine.\_pretrain\_metrics} and
\py{eval.\_accumulate}.

% ---------------------------------------------------------------------
\subsection{For \texorpdfstring{\cref{sec:se2-target}}{Section 3.1} --- the \texorpdfstring{\SEtwo}{SE(2)} target}
% ---------------------------------------------------------------------

\paragraph{Patch state as a group element.} Sampling gives patch $i$ a centre
$c_i \in \mathbb{R}^2$ and an orientation $\varphi_i$, i.e. an element of \SEtwo{}:
\begin{equation}\label{eq:se2-element}
  g_i \;=\; (c_i, \varphi_i) \;\cong\;
  \begin{bmatrix} R(\varphi_i) & c_i \\ 0^{\!\top} & 1 \end{bmatrix} \in \SEtwo .
\end{equation}
PART's patches are axis-aligned, so $\varphi_i \equiv 0$ and $g_i$ lives in the
translation subgroup $\mathbb{R}^2 \triangleleft \SEtwo$. This is the one-line
statement of what the project extends.

\paragraph{The relative transform.} The group-theoretically correct relative
transform between the ordered pair is
\begin{equation}\label{eq:se2-relative}
  g_i^{-1} g_j \;=\; \bigl(\, R(-\varphi_i)(c_j - c_i),\; \varphi_j - \varphi_i \,\bigr).
\end{equation}

\paragraph{What is actually supervised.} \textbf{RoPART does not regress
\cref{eq:se2-relative}.} The translation channels are differenced in the
\emph{global} image frame, not rotated into the reference patch's frame
(\py{build\_targets}); only the orientation channel is genuinely relative.
The implemented target for the ordered pair $(i,j)$ is
\begin{equation}\label{eq:target-ropart}
  \theta_{ij}
  \;=\;
  \begin{bmatrix}
    \Delta x \\ \Delta y \\ \cos\Delta\varphi \\ \sin\Delta\varphi
  \end{bmatrix}_{ij}
  \;=\;
  \begin{bmatrix}
    (c^{x}_{j} - c^{x}_{i}) / P \\[2pt]
    (c^{y}_{j} - c^{y}_{i}) / P \\[2pt]
    \cos(\varphi_j - \varphi_i) \\[2pt]
    \sin(\varphi_j - \varphi_i)
  \end{bmatrix},
  \qquad
  \Delta\varphi \;\triangleq\; \varphi_j - \varphi_i .
\end{equation}
The two frames coincide up to the $R(-\varphi_i)$ factor; the reference-local variant
of \cref{eq:se2-relative} is the deferred ablation recorded in
\cref{sec:position-convention}. \emph{Write this honestly: the target is
``\SEtwo{} with the translation read in a global frame'', not $g_i^{-1}g_j$.}

Note also that $P$ in \cref{eq:target-ropart} is a constant, whereas PART divides by
the reference patch's own $w_{\mathrm{ref}}, h_{\mathrm{ref}}$
(\cref{sec:part}). RoPART fixes every patch to $P \times P$, so the two normalisations
are identical here --- but say so, because it is also what removes the
$(\Delta w, \Delta h)$ signal.

\paragraph{The dense table.} For a batch of $b$ images the builder emits every ordered
pair at once,
\begin{equation}\label{eq:target-table}
  T \in \mathbb{R}^{b \times C \times N \times N},
  \qquad T[\,\cdot\,, :, i, j] = \theta_{ij},
\end{equation}
with the build order fixing $C$:
\begin{equation}\label{eq:num-channels}
  C \;=\;
  \begin{cases}
    2, & (\Delta x, \Delta y) \quad \text{--- translation-only baseline},\\
    4, & +\,(\cos\Delta\varphi, \sin\Delta\varphi) \quad \text{--- RoPART},\\
    6, & +\,(\Delta w, \Delta h) \quad \text{--- deferred, \textbf{never run}}.
  \end{cases}
\end{equation}
$C$ is also the output width of the relative head (\cref{eq:head-out}), which is
what makes the build order a one-line change.

\paragraph{Structural invariants.} These are asserted on the \emph{targets} at
construction (\py{check\_*\_invariants}) and \emph{measured} on the
predictions in \cref{sec:protocol-and-metrics}:
\begin{align}
  \theta^{xy}_{ii} &= 0, &
  (\cos\Delta\varphi_{ii}, \sin\Delta\varphi_{ii}) &= (1, 0),
  \label{eq:inv-identity}\\
  \theta^{xy}_{ij} &= -\,\theta^{xy}_{ji}, &
  \cos\Delta\varphi_{ij} &= \phantom{-}\cos\Delta\varphi_{ji},
  \quad
  \sin\Delta\varphi_{ij} = -\sin\Delta\varphi_{ji},
  \label{eq:inv-negsym}\\
  \theta^{xy}_{ik} &= \theta^{xy}_{ij} + \theta^{xy}_{jk}, &
  \Delta\varphi_{ik} &= \operatorname{wrap}(\Delta\varphi_{ij} + \Delta\varphi_{jk}).
  \label{eq:inv-composition}
\end{align}
\Cref{eq:inv-negsym} is the rotational analogue of PART's Figure~6; the asymmetry of
the two rotation columns (cosine even, sine odd) is the sharper statement, since a
head that merely predicted a constant would satisfy the cosine half.

% ---------------------------------------------------------------------
\subsection{For \texorpdfstring{\cref{sec:angle-representation}}{Section 3.2} --- angle representation and loss}
% ---------------------------------------------------------------------

\paragraph{Why not a raw angle.} A scalar parametrisation $\varphi \in (-\pi,\pi]$ makes
the squared error discontinuous across the seam: for $\varepsilon \to 0^{+}$,
\begin{equation}\label{eq:wrap-discontinuity}
  \bigl(\,(\pi - \varepsilon) - (-\pi + \varepsilon)\,\bigr)^2 \;\longrightarrow\; 4\pi^2,
  \qquad\text{while the true angular separation is } 2\varepsilon \to 0 .
\end{equation}
This is the concrete form of the objection in \cite{rotcont}, and it is also
\textbf{Hard-won fact \#3}: \py{WeightedMSELoss} has an all-weights-one fast path that
would apply exactly \cref{eq:wrap-discontinuity} to a raw angle channel.

\paragraph{The circle encoding.} Regress the point on $S^1$ instead,
\begin{equation}\label{eq:circle-encoding}
  \Delta\varphi \;\longmapsto\; (\cos\Delta\varphi,\, \sin\Delta\varphi) \in S^1,
\end{equation}
and recover the predicted angle from the (unconstrained) head outputs
$(\hat c, \hat s)$ by
\begin{equation}\label{eq:atan2}
  \widehat{\Delta\varphi} \;=\; \operatorname{atan2}(\hat s, \hat c).
\end{equation}

\paragraph{MSE on the circle \emph{is} a cosine loss.} If the prediction happened to lie
on the unit circle, the two candidate losses named in the text coincide up to a factor
of two:
\begin{equation}\label{eq:mse-equals-cosine}
  (\hat c - \cos\Delta\varphi)^2 + (\hat s - \sin\Delta\varphi)^2
  \;=\; \lVert \hat u \rVert^2 + 1 - 2\,\hat u^{\!\top} u
  \;\overset{\lVert\hat u\rVert = 1}{=}\; 2\bigl(1 - \cos(\widehat{\Delta\varphi} - \Delta\varphi)\bigr),
\end{equation}
with $\hat u = (\hat c, \hat s)$ and $u = (\cos\Delta\varphi, \sin\Delta\varphi)$.
\emph{Caveat to state:} the head output is \emph{not} norm-constrained, so
\cref{eq:mse-equals-cosine} holds only up to the $\lVert\hat u\rVert^2$ term ---
plain MSE additionally penalises radial deviation, which acts as a free confidence
channel. The runs used MSE, not the cosine form.

\paragraph{The loss as run.} Grouped MSE over $M$ sampled pairs, with the two channel
groups $\mathcal{G}_{xy} = \{0,1\}$ and $\mathcal{G}_{\varphi} = \{2,3\}$
(\py{losses.RelativeMSE}). The unbalanced form --- upstream's, and the default that
keeps the baseline path unchanged --- is
\begin{equation}\label{eq:loss-none}
  \mathcal{L}
  \;=\; w_{xy}\,\frac{1}{2 b M}\!\!\sum_{k \in \mathcal{G}_{xy}}\sum_{n,m}\!\bigl(\hat\theta_{knm} - \theta_{knm}\bigr)^2
  \;+\; w_{\varphi}\,\frac{1}{2 b M}\!\!\sum_{k \in \mathcal{G}_{\varphi}}\sum_{n,m}\!\bigl(\hat\theta_{knm} - \theta_{knm}\bigr)^2 .
\end{equation}
Setting $w_\varphi = 0$ and dropping $\mathcal{G}_\varphi$ recovers PART's
$\mathcal{L}_{\mathrm{MSE}}$ of \cref{sec:part} exactly.

\paragraph{Variance balancing --- the form every reported run actually used.} Each
channel's squared error is divided by that channel's own detached batch target
variance, i.e.\ by its own mean-predictor floor:
\begin{equation}\label{eq:loss-variance}
  \mathcal{L}
  \;=\;
  \sum_{G \in \{\mathcal{G}_{xy}, \mathcal{G}_{\varphi}\}} \!\! w_G \,
  \frac{1}{|G|} \sum_{k \in G}
  \frac{\operatorname{MSE}_k}{\max\!\bigl(\operatorname{sg}[\operatorname{Var}(\theta_k)],\, \epsilon\bigr)},
  \qquad \epsilon = 10^{-3},
\end{equation}
where $\operatorname{sg}[\cdot]$ is the stop-gradient and both $\operatorname{MSE}_k$ and
$\operatorname{Var}(\theta_k)$ are taken over the pooled $(b \times M)$ batch sample
(biased variance). Every term is then a \emph{fraction of its own floor}, so both
groups are $O(1)$ irrespective of units or rotation range, and $w_{xy}, w_\varphi$
become genuine relative-importance knobs.

\emph{Two points the prose must make, because both are load-bearing:}
\begin{enumerate}
  \item \textbf{The need for \cref{eq:loss-variance} is quantitative, not aesthetic.}
  With \cref{eq:loss-none} and $w_{xy}=w_\varphi=1$ the translation term is
  $\operatorname{Var}(\Delta x) \approx 6$ against a rotation group bounded by
  $1 - (\sin a / a)^4$ (\cref{eq:rot-floor-bounded}), which at $a = 30^\circ$ is
  $0.169$ --- a $\sim$\,$30\times$ imbalance, and $\Delta\varphi$ collapses to the mean
  predictor.
  \item \textbf{$\epsilon$ is not a formality.} As $a \to 0$,
  $\operatorname{Var}(\cos\Delta\varphi) = O(a^4)$, so an unclamped divisor amplifies
  that channel by $\sim 10^{5}$--$10^{6}$ and NaNs training; $\epsilon = 10^{-3}$ caps
  the amplification at $\sim 10^{3}$ and is inert for the $\pm 30^\circ$, $\pm 90^\circ$
  and quad runs, whose variances are $\gg \epsilon$.
\end{enumerate}
The reported arms use $w_{xy} = 1$, $w_\varphi = 0.1$, and the translation-only baseline
is run under \cref{eq:loss-variance} too --- otherwise its effective gradient scale
differs from RoPART's translation group by $\sim$\,$20\times$ and it is not a control.

\paragraph{\textbf{NOT IMPLEMENTED} --- Kendall homoscedastic weighting.} Cited in the
prose as the established adaptive alternative \cite{kendall}; write it only as
related work, never as something run:
\begin{equation}\label{eq:kendall}
  \mathcal{L} \;=\; \sum_{G} \frac{1}{2\sigma_G^{2}} \mathcal{L}_G \;+\; \log \sigma_G ,
  \qquad \sigma_G \text{ learned}.
\end{equation}
\Cref{eq:loss-variance} is the fixed, data-driven analogue: it divides by an
\emph{observed} target variance rather than a learned uncertainty.

% ---------------------------------------------------------------------
\subsection{For \texorpdfstring{\cref{sec:position-convention}}{Section 3.3} --- patch position convention}
% ---------------------------------------------------------------------

\paragraph{The two anchors.} For the axis-aligned box $[x_s, y_s, x_e, y_e]$ with
$x_e = x_s + P$, $y_e = y_s + P$:
\begin{equation}\label{eq:anchors}
  c_i^{\text{centre}} = \Bigl(\tfrac{x_s + x_e}{2},\, \tfrac{y_s + y_e}{2}\Bigr)
  \qquad\text{vs.}\qquad
  c_i^{\text{corner}} = (x_s, y_s).
\end{equation}

\paragraph{Why the choice is free at $\varphi = 0$ but not after.} With constant patch
size the two anchors differ by the same fixed offset for every patch, so the offset
cancels in the difference:
\begin{equation}\label{eq:anchor-equivalence}
  c^{\text{centre}}_j - c^{\text{centre}}_i
  \;=\; \bigl(c^{\text{corner}}_j + \tfrac{P}{2}\mathbf{1}\bigr) - \bigl(c^{\text{corner}}_i + \tfrac{P}{2}\mathbf{1}\bigr)
  \;=\; c^{\text{corner}}_j - c^{\text{corner}}_i .
\end{equation}
Hence \emph{no baseline number changes} under the convention switch. Once patches carry
$\varphi_i$, the corner is no longer a fixed offset from the centre --- it moves with
the patch:
\begin{equation}\label{eq:corner-entangles}
  c^{\text{corner}}_i(\varphi_i) \;=\; c^{\text{centre}}_i - R(\varphi_i)\,\tfrac{P}{2}\mathbf{1},
\end{equation}
so a corner-anchored $\Delta x, \Delta y$ would carry a $\varphi$-dependent term and the
translation channel would no longer be a clean control for the rotation channel. The
centre is the unique fixed point of $R(\varphi_i)$ about the patch, which is the whole
argument.

\paragraph{The half-pixel subtlety --- worth a footnote, not a paragraph.} The target
builder uses the geometric centre $(x_s + x_e)/2$, while the rotated pixel read uses the
\emph{pixel-index} centre $(x_s + x_e - 1)/2$ so that an integer coordinate lands exactly
on a pixel centre under \py{align\_corners=True} (which is what keeps
$\varphi \in \{0, 90^\circ, \dots\}$ interpolation-free). The two differ by a constant
$\tfrac{1}{2}\mathbf{1}$ for every patch, so by the same cancellation as
\cref{eq:anchor-equivalence} the targets are unaffected.

\paragraph{\textbf{NOT IMPLEMENTED} --- the reference-local frame.} The deferred
alternative, and the translation half of \cref{eq:se2-relative}:
\begin{equation}\label{eq:local-frame}
  \Delta p^{\,\text{local}}_{ij} \;=\; \tfrac{1}{P}\, R(-\varphi_i)\,(c_j - c_i).
\end{equation}

% ---------------------------------------------------------------------
\subsection{For \texorpdfstring{\cref{sec:rotating-pixels}}{Section 3.4} --- rotating pixels}
% ---------------------------------------------------------------------

\paragraph{The empty-corner problem, and the window that solves it.} A $P \times P$
square rotated by an arbitrary angle is not contained in its own axis-aligned footprint;
it is contained in its diagonal. So read from
\begin{equation}\label{eq:rot-window}
  W \;=\; \bigl\lceil P\sqrt{2} \,\bigr\rceil,
  \qquad
  m \;=\; \Bigl\lfloor \tfrac{W - P + 1}{2} \Bigr\rfloor ,
\end{equation}
where $m$ is the edge margin a patch centre needs for the diagonal window to stay
in-image. At $P = 32$: $W = 46$, $m = 7$. \emph{This margin is not cosmetic --- it
narrows the sampling range and therefore changes the translation floor
(\cref{eq:trans-floor}), which is why the baseline and rotation arms are not scored
against an identical target distribution.}

\paragraph{Box sampling.} Top-left corners are drawn independently and uniformly over
the valid integer offsets,
\begin{equation}\label{eq:box-sampling}
  x_s, y_s \;\sim\; \mathcal{U}\{\,m,\, m+1,\, \dots,\, S - P - m\,\},
  \qquad
  M_{\text{pos}} \;=\; S - P - 2m + 1 ,
\end{equation}
with $m = 0$ for the translation baseline and for the quad control (exact
\py{rot90} needs no diagonal window). Positions are continuous in the off-grid sense
--- any pixel offset, not a multiple of $P$ --- and patches may overlap.

\paragraph{Angle sampling.} One orientation per patch, i.i.d.:
\begin{equation}\label{eq:angle-sampling}
  \varphi_i \;\sim\;
  \begin{cases}
    \mathcal{U}[0, 2\pi) & \text{full circle},\\
    \mathcal{U}[-a, +a] & \text{bounded, } a \in \{30^\circ, 60^\circ, 90^\circ\},\\
    \mathcal{U}\{0, \tfrac{\pi}{2}, \pi, \tfrac{3\pi}{2}\} & \text{quad}.
  \end{cases}
\end{equation}

\paragraph{The read itself.} Extraction is a single inverse warp: for output pixel at
patch-local, centred coordinates $(l_x, l_y)$ with
$l \in \{-\tfrac{P-1}{2}, \dots, \tfrac{P-1}{2}\}$, the source coordinate is
\begin{equation}\label{eq:inverse-warp}
  \begin{bmatrix} s_x \\ s_y \end{bmatrix}
  \;=\;
  \begin{bmatrix} c^{x}_i \\ c^{y}_i \end{bmatrix}
  + R(\varphi_i)
  \begin{bmatrix} l_x \\ l_y \end{bmatrix},
\end{equation}
and the patch value is the bilinear sample $p(l_x, l_y) = I(s_x, s_y)$. \emph{Exactly one
interpolation} --- crop--rotate--centre-crop--resize is collapsed into one
\py{grid\_sample}, which matters because each additional resampling would compound the
blur of \cref{sec:interpolation-controls}. The normalisation into
\py{grid\_sample} coordinates is against the \emph{original} extent, with
\py{align\_corners=True}:
\begin{equation}\label{eq:grid-normalise}
  g_x \;=\; \frac{s_x}{(S-1)/2} - 1, \qquad g_y \;=\; \frac{s_y}{(S-1)/2} - 1 ,
\end{equation}
so that a supersampled source (\cref{eq:supersample}) maps to the same physical
coordinates without rescaling the grid.

\paragraph{Bilinear interpolation, written out.} Needed only if
\cref{sec:interpolation-controls} argues about the point-spread function; with
$\delta = s - \lfloor s \rfloor$ the fractional part,
\begin{equation}\label{eq:bilinear}
  I(s_x, s_y) \;=\; \sum_{u,v \in \{0,1\}}
    \bigl(1 - |u - \delta_x|\bigr)\bigl(1 - |v - \delta_y|\bigr)\;
    I\bigl[\lfloor s_x \rfloor + u,\; \lfloor s_y \rfloor + v\bigr].
\end{equation}
The kernel is separable with support that depends on $(\delta_x, \delta_y)$, and
$(\delta_x, \delta_y)$ is a function of $\varphi_i$ --- that dependence \emph{is} the
confound.

\paragraph{The quad path is a permutation, not a resampling.} With
\begin{equation}\label{eq:quad-k}
  k \;=\; \Bigl(\operatorname{round}\!\bigl(\varphi_i \big/ \tfrac{\pi}{2}\bigr)\Bigr) \bmod 4,
  \qquad
  p \;=\; \operatorname{rot90}^{k}\bigl(\text{crop}_{P\times P}\bigr),
\end{equation}
$\delta_x = \delta_y = 0$ identically, \cref{eq:bilinear} degenerates to a copy, and no
angle-correlated blur exists. This is the sense in which quad is confound-free
\emph{by construction} rather than by control.

\paragraph{Packing into the pseudo-image.} Patch $k$ is placed row-major in a
$\sqrt{N} \times \sqrt{N}$ tiling, so the ViT's stride-$P$ convolution recovers the
tokens one-to-one:
\begin{equation}\label{eq:retile}
  \hat I\bigl[\,\cdot\,,\; r P + u,\; c P + v\,\bigr] \;=\; p_k[\,\cdot\,, u, v],
  \qquad r = \bigl\lfloor k / \sqrt{N} \bigr\rfloor,\; c = k \bmod \sqrt{N}.
\end{equation}
\emph{Say plainly that the tiling order carries no geometric information} --- $k$ is
the sampling index, unrelated to where the patch came from. That is what makes the
pseudo-image an unordered set to a position-embedding-free encoder.

Finally the standard channel normalisation, applied after packing:
\begin{equation}\label{eq:normalise}
  \tilde I \;=\; (\hat I - \mu)/\sigma,
  \quad \mu = (0.485, 0.456, 0.406),
  \quad \sigma = (0.229, 0.224, 0.225).
\end{equation}
Worth one sentence: pretraining applies \emph{no other} augmentation, deliberately, so
that the baseline-vs-rotation contrast is not confounded by the augmentation pipeline
(upstream applies full DeiT augmentation).

% ---------------------------------------------------------------------
\subsection{For \texorpdfstring{\cref{sec:interpolation-controls}}{Section 3.5} --- interpolation controls}
% ---------------------------------------------------------------------

\paragraph{The signature, made measurable.} The confound needs a scalar before it can be
controlled. Use mean absolute pixel gradient --- a high-frequency-energy proxy --- of a
patch rotated by $\varphi$:
\begin{equation}\label{eq:sharpness}
  E(\varphi) \;=\; \bigl\langle\, |\partial_x p_\varphi| + |\partial_y p_\varphi| \,\bigr\rangle .
\end{equation}
Sweeping $\varphi$ gives a curve that peaks at $0^\circ, 90^\circ, 180^\circ$ and dips at
$45^\circ, 135^\circ$ --- so $E$ is a deterministic, angle-correlated leak. This is the
sentence that turns ``interpolation noise'' into ``a false-positive generator''.

\paragraph{The trap: the blur is oriented.} Bilinear resampling is convolution with the
separable kernel of \cref{eq:bilinear}, whose effective point-spread function
$\mathrm{PSF}_{\text{interp}}(\varphi)$ is anisotropic with an orientation that tracks
$\varphi$. Therefore
\begin{equation}\label{eq:magnitude-not-sufficient}
  E_{\text{ctrl}}(\varphi) = \text{const}
  \;\;\not\Longrightarrow\;\;
  \mathrm{PSF}_{\text{ctrl}}(\varphi) = \text{const} .
\end{equation}
\emph{This is the key inequality of the section:} flattening the magnitude cue is
necessary but not sufficient, and it is what ranks the four controls below.

\paragraph{(a) Matched blur --- \textbf{NOT IMPLEMENTED}.} Solve for a per-angle
isotropic $\sigma(\varphi)$ that flattens \cref{eq:sharpness} to its worst case:
\begin{equation}\label{eq:matched-blur}
  \sigma(\varphi) \;:\; E\bigl(G_{\sigma(\varphi)} * p_\varphi\bigr) \;=\; \min_{\varphi'} E(\varphi')
  \quad \forall \varphi .
\end{equation}
Prototyped by bisection ($E$ is monotone decreasing in $\sigma$) but never wired into a
run --- \py{\_extract\_matched} raises \py{NotImplementedError}. By
\cref{eq:magnitude-not-sufficient} it addresses only the magnitude cue anyway.

\paragraph{(b) Randomised blur.} Draw the blur independently of the angle,
\begin{equation}\label{eq:randomised-blur}
  \sigma \sim \mathcal{U}(0, \sigma_{\max}), \quad \sigma_{\max} = 0.8,
  \qquad p \;\leftarrow\; G_\sigma * p_\varphi .
\end{equation}
State its limitation as an equation, because it is counter-intuitive: randomisation adds
angle-independent \emph{variance} but leaves the systematic term intact,
\begin{equation}\label{eq:randomised-limit}
  \mathbb{E}\bigl[E \mid \varphi\bigr] \;\text{still dips at } 45^\circ .
\end{equation}
So (b) is the \emph{weakest} control, not the cheap win it appears to be.

\paragraph{(c) Supersampling --- the control the results rest on.} Bicubically upsample
the source by $k$ before the read of \cref{eq:inverse-warp}, so the fractional offsets
of \cref{eq:bilinear} shrink:
\begin{equation}\label{eq:supersample}
  \delta^{(k)} \;=\; O(1/k), \qquad k = 4 ,
\end{equation}
making the resampling near-lossless and its PSF near angle-independent --- it removes
both the magnitude \emph{and} the direction cue. Empirically $k = 4$ leaves a residual
sharpness ratio of $0.96$ against $0.81$ for the raw path, i.e.\ $\approx 96\%$ of the
signature removed, not $100\%$. \emph{Quote that residual; it is the honest loose end of
the whole positive result.}

\paragraph{(d) Dominant isotropic low-pass.} Swamp the anisotropy with one heavy
isotropic blur on every patch:
\begin{equation}\label{eq:dominant-lp}
  \mathrm{PSF}_{\text{eff}}(\varphi) \;=\; \mathrm{PSF}_{\text{interp}}(\varphi) * G_{\sigma_{\mathrm{lp}}},
  \qquad \sigma_{\mathrm{lp}} = 1.0 ,
\end{equation}
chosen so the isotropic term dominates and the $\varphi$-dependence becomes a
perturbation. Heaviest signal loss of the four.

\paragraph{Discretised Gaussian.} If the implementation detail is wanted: the kernel side
is $\max\bigl(2\operatorname{round}(3\sigma) + 1,\, 3\bigr)$, i.e.\ $\pm 3\sigma$
support forced odd.

\paragraph{The control matrix.} The point of the section in one line --- the controls
differ in \emph{which} cue they break:
\begin{equation}\label{eq:control-summary}
  \text{(a) magnitude only},\quad
  \text{(b) neither, adds variance},\quad
  \text{(c) both},\quad
  \text{(d) both, at heavy cost}.
\end{equation}

% ---------------------------------------------------------------------
\subsection{For \texorpdfstring{\cref{sec:relative-head}}{Section 3.6} --- the relative head}
% ---------------------------------------------------------------------

\paragraph{Encoder.} A plain pre-norm ViT, with the one departure that matters:
\begin{align}
  X &= \mathrm{PatchEmbed}(\tilde I) \in \mathbb{R}^{N \times d},
  &
  Z^{(0)} &= \bigl[\,X \,;\, x_{\mathrm{cls}}\,\bigr] \in \mathbb{R}^{(N+1) \times d},
  \label{eq:vit-embed}\\
  \tilde Z^{(\ell)} &= Z^{(\ell-1)} + \mathrm{MHSA}\bigl(\mathrm{LN}(Z^{(\ell-1)})\bigr),
  &
  Z^{(\ell)} &= \tilde Z^{(\ell)} + \mathrm{MLP}\bigl(\mathrm{LN}(\tilde Z^{(\ell)})\bigr),
  \label{eq:vit-block}\\
  X' &= \mathrm{LN}\bigl(Z^{(L)}\bigr)_{1:N} \in \mathbb{R}^{N \times d}.
  &&
  \label{eq:vit-out}
\end{align}
\textbf{No position embeddings are added during pretraining} --- $E_{\mathrm{pos}}$ is
present as a parameter but is applied only on the classification forward. This is
PART's premise and RoPART inherits it: with attention permutation-equivariant and
\cref{eq:retile} carrying no geometry, the encoder sees an unordered set, so the
geometry \emph{is} the supervision. Note the \py{[CLS]} token is appended \emph{last},
not prepended, and is dropped before the relative head --- worth a footnote because
it is the reason a frozen linear probe on this checkpoint is ill-posed.

\paragraph{Pair sampling.} $M$ ordered pairs are drawn without replacement from the
$N^2$ cells of the table, once per forward and \emph{shared across the batch}:
\begin{equation}\label{eq:pair-sampling}
  \mathrm{idx} \subset \{0, \dots, N^2 - 1\},\; |\mathrm{idx}| = M,
  \qquad
  i = \bigl\lfloor \mathrm{idx} / N \bigr\rfloor,
  \quad
  j = \mathrm{idx} \bmod N .
\end{equation}
Give the numbers: at $P = 32$, $N = 49$ and $M = 2048$ of $N^2 = 2401$ --- $85\%$ of all
ordered pairs, including the $N$ diagonal cells, so the identity constraint
\cref{eq:inv-identity} is in-distribution rather than an extrapolation. Contrast with
the $M = 64$ default we inherited; upstream's own sweep uses $512$--$4096$.

\paragraph{The query --- the decision of \cref{sec:relative-head}.} For pair
$m = (i, j)$, the two variants are
\begin{align}
  q_m^{\texttt{patch\_cat}} &= W_{\mathrm{in}}\bigl[\,x'_i \,;\, x'_j\,\bigr] + b_{\mathrm{in}},
  \qquad W_{\mathrm{in}} \in \mathbb{R}^{d \times 2d},
  \label{eq:query-patchcat}\\
  q_m^{\texttt{positional}} &= \mathrm{PE}_{2\mathrm{d}}[i, j, :],
  \qquad\text{with keys } x'_n \leftarrow x'_n + \mathrm{PE}_{1\mathrm{d}}[n].
  \label{eq:query-positional}
\end{align}
\emph{The contrast is the point.} \Cref{eq:query-patchcat} builds the query from
content alone, so no positional code enters the head at any point and the ``no position
embeddings'' premise survives end to end; \cref{eq:query-positional} injects a fixed
sinusoidal code for the pair index into both query and keys, and on this configuration
the baseline then collapses to the mean-predictor floor and cannot overfit a single
batch. \Cref{eq:query-patchcat} is the default for baseline and RoPART alike.

\paragraph{Cross-attention.} One \py{nn.MultiheadAttention} layer with a single head,
queries $Q = [q_1; \dots; q_M]$ and keys/values the patch tokens $X'$:
\begin{equation}\label{eq:cross-attention}
  Z \;=\; \operatorname{softmax}\!\left(\frac{(Q W_Q)(X' W_K)^{\!\top}}{\sqrt{d}}\right)(X' W_V)\, W_O .
\end{equation}
\emph{Note the divergence from \cref{sec:part}:} the paper writes the head as
$\operatorname{softmax}(\hat X X'^{\!\top}/\sqrt{d})\,X'$, omitting the four internal
projections $W_Q, W_K, W_V, W_O$ that \py{nn.MultiheadAttention} in fact applies. Either
present \cref{eq:cross-attention} and say the paper's form is the schematic, or present
the schematic and footnote the projections --- but do not silently mix them.

\paragraph{Output.} A final linear map to the channel width of
\cref{eq:num-channels}:
\begin{equation}\label{eq:head-out}
  \hat\theta \;=\; Z\,W_{\mathrm{out}} + b_{\mathrm{out}} \in \mathbb{R}^{M \times C},
  \qquad W_{\mathrm{out}} \in \mathbb{R}^{d \times C}.
\end{equation}
$W_{\mathrm{out}}$ is the \emph{only} tensor whose shape depends on the build order ---
which is the concrete content of ``the head is fully parametrised by
\py{num\_channels}''.

% ---------------------------------------------------------------------
\subsection{For \texorpdfstring{\cref{sec:implementation}}{Section 3.7} and \texorpdfstring{\cref{sec:infrastructure}}{Section 3.8}}
% ---------------------------------------------------------------------

\stub{No equations needed. The one formal statement worth making in
\cref{sec:implementation} is the flag-gating invariant --- that with rotation off the
computation is identical to the baseline path, i.e.\
$\text{RoPART}\big|_{C=2} \equiv \text{baseline}$ --- and that is better made as a
sentence plus the test name than as display maths.}

% ---------------------------------------------------------------------
\subsection{For \texorpdfstring{\cref{sec:protocol-and-metrics}}{Section 3.9} --- protocol and metrics}
% ---------------------------------------------------------------------

\paragraph{Optimisation.} AdamW, $\beta = (0.9, 0.999)$, weight decay $0.05$ applied to
matrices only (biases, norms, \py{cls\_token} and \py{pos\_embed} excluded). The
schedule is per-\emph{epoch}, linear warmup into cosine decay:
\begin{equation}\label{eq:lr-schedule}
  \eta(e) \;=\;
  \begin{cases}
    \eta_{\mathrm{warm}} + (\eta_{\max} - \eta_{\mathrm{warm}})\dfrac{e}{E_w},
      & e < E_w, \\[10pt]
    \eta_{\min} + (\eta_{\max} - \eta_{\min})\,\dfrac{1}{2}\!\left(1 + \cos \pi \dfrac{e - E_w}{E - E_w}\right),
      & e \geq E_w,
  \end{cases}
\end{equation}
with $\eta_{\max} = 5\!\times\!10^{-4}$, $\eta_{\mathrm{warm}} = 10^{-6}$,
$\eta_{\min} = 10^{-5}$, $E_w = 5$, $E = 250$, batch size $192$.
\emph{Two things to state:} (i) unlike upstream's \py{main.py}, no linear
$\eta \times \text{batch}/512$ rescaling is applied --- $\eta_{\max}$ is the literal
peak; (ii) $\eta$ is a pure function of $e$ and $E$, which is why a longer run had to be
launched from scratch rather than resumed (resuming with a larger $E$ would warm-restart
the cosine near its peak).

\paragraph{Mean-predictor floors --- derive these once, here.} Every result in
\cref{ch:pretraining} is quoted as a fraction of a floor, so the floors must be
exact rather than asserted.

\emph{Translation.} With corners i.i.d.\ uniform over $M_{\text{pos}}$ integer offsets
(\cref{eq:box-sampling}) and $\Delta x = (x_{s,j} - x_{s,i})/P$,
\begin{equation}\label{eq:trans-floor}
  \operatorname{Var}(\Delta x)
  \;=\; \frac{2\operatorname{Var}(x_s)}{P^2}
  \;=\; \frac{M_{\text{pos}}^{2} - 1}{6 P^{2}} .
\end{equation}
At $S = 224$, $P = 32$: the baseline ($m = 0$, $M_{\text{pos}} = 193$) gives $6.06$ and
the rotation arms ($m = 7$, $M_{\text{pos}} = 179$) give $5.21$, both matching the
measured \py{floor\_x} to $0.1\%$. \emph{This equation is what licenses the caveat that
the arms are not measured against an identical target distribution}, and the
rescaling factor $\sqrt{6.06/5.21} = 1.079$ used to check that the ordering does not
depend on the choice.

\emph{Rotation, full circle.} Predicting $(\hat c, \hat s) = (0,0)$ gives
\begin{equation}\label{eq:rot-floor-uniform}
  \mathbb{E}\bigl[\cos^2\Delta\varphi + \sin^2\Delta\varphi\bigr] \;=\; 1 \quad\text{exactly},
\end{equation}
the origin of the ``$\text{rot\_mse}$ floor $= 1.0$'' quoted throughout.

\emph{Rotation, bounded.} For $\varphi \sim \mathcal{U}[-a, a]$ i.i.d.,
$\mathbb{E}[\cos\varphi] = \sin a / a$ and $\mathbb{E}[\sin\varphi] = 0$, so
$\mathbb{E}[\cos\Delta\varphi] = (\sin a / a)^2$ and the floor is
\begin{equation}\label{eq:rot-floor-bounded}
  \operatorname{Var}(\cos\Delta\varphi) + \operatorname{Var}(\sin\Delta\varphi)
  \;=\; 1 - \Bigl(\frac{\sin a}{a}\Bigr)^{\!4}
  \;=\;
  \begin{cases} 0.169, & a = 30^\circ,\\ 0.836, & a = 90^\circ.\end{cases}
\end{equation}
\textbf{Flag this in the text:} the training log hard-codes $\text{rot\_floor} = 1.0$,
which is exact only for the full circle, so the logged bounded-run ratio is
conservative. The reported bounded results are quoted against the \emph{angular} floor
below, which is computed empirically and is correct.

\emph{Angular floor.} $\Delta\varphi = \varphi_j - \varphi_i$ with $\varphi \sim
\mathcal{U}[-a,a]$ is triangular on $[-2a, 2a]$ with density
$f(t) = (2a - |t|)/(4a^2)$, so the constant predictor $\Delta\varphi = 0$ --- which is
the circular median --- achieves
\begin{equation}\label{eq:ang-floor}
  \mathbb{E}\bigl|\Delta\varphi\bigr| \;=\; \tfrac{2}{3} a
  \;=\;
  \begin{cases} 20.0^\circ, & a = 30^\circ,\\ 60.0^\circ, & a = 90^\circ,\end{cases}
  \qquad\text{and } 90^\circ \text{ for the full circle.}
\end{equation}
Measured: $19.96^\circ$ and $59.89^\circ$. \emph{This is the most useful equation in the
chapter} --- it is the number every rotation result is read against, and it explains why
$\pm 90^\circ$ can have a worse absolute angular error and still be learning
orientation \emph{better} once floor-normalised.

\paragraph{Intrinsic gates.} All four are computed on the \emph{dense} $N \times N$
prediction table (every ordered pair, not the training subsample), off-diagonal except
where stated. Write $\hat t_{ij} = (\hat\theta^x_{ij}, \hat\theta^y_{ij})$ and
$\hat u_{ij} = (\hat c_{ij}, \hat s_{ij})$.

\emph{(1) Floor ratios and pixel error.}
\begin{equation}\label{eq:metric-floor-ratio}
  \frac{\operatorname{MSE}_x}{\operatorname{Var}(\Delta x)},
  \qquad
  \mathrm{RMSE}_{\mathrm{px}} \;=\; P\sqrt{\bigl\langle\, \lVert \hat t_{ij} - t_{ij} \rVert^2 \,\bigr\rangle_{i \neq j}} .
\end{equation}

\emph{(2) Identity.} Query the head on $(i,i)$, where the answer is known
(\cref{eq:inv-identity}):
\begin{equation}\label{eq:metric-identity}
  \mathrm{Id}_{xy} = P\sqrt{\bigl\langle \lVert \hat t_{ii} \rVert^2 \bigr\rangle_i},
  \qquad
  \mathrm{Id}_{\mathrm{rot}} = \sqrt{\bigl\langle (\hat c_{ii} - 1)^2 + \hat s_{ii}^2 \bigr\rangle_i} .
\end{equation}
Its diagnostic value is worth a sentence: a head that ignored the queried slots and
averaged over the image would fail \cref{eq:metric-identity} while still beating the
floor.

\emph{(3) Negative symmetry}, as a residual-to-signal RMS ratio so it is scale-free:
\begin{align}
  \rho^{\mathrm{sym}}_{xy}
    &= \sqrt{\frac{\sum_{i \neq j} \lVert \hat t_{ij} + \hat t_{ji} \rVert^2}
                  {\sum_{i \neq j} \lVert \hat t_{ij} \rVert^2}},
  \label{eq:metric-negsym-xy}\\
  \rho^{\mathrm{sym}}_{\mathrm{rot}}
    &= \sqrt{\frac{\sum_{i \neq j} \bigl[(\hat c_{ij} - \hat c_{ji})^2 + (\hat s_{ij} + \hat s_{ji})^2\bigr]}
                  {\sum_{i \neq j} \bigl[\hat c_{ij}^2 + \hat s_{ij}^2\bigr]}} .
  \label{eq:metric-negsym-rot}
\end{align}
Both $\to 0$ for a perfectly consistent head. \Cref{eq:metric-negsym-rot} tests the
cosine-even/sine-odd pair of \cref{eq:inv-negsym} simultaneously.

\emph{(4) Compositional consistency}, over random distinct triples $(i,j,k)$ --- the
strongest gate, since pairwise symmetry can hold while composition fails:
\begin{equation}\label{eq:metric-comp-xy}
  \rho^{\mathrm{comp}}_{xy}
  = \sqrt{\frac{\sum \lVert \hat t_{ik} - (\hat t_{ij} + \hat t_{jk}) \rVert^2}
               {\sum \lVert \hat t_{ik} \rVert^2}},
\end{equation}
\begin{equation}\label{eq:metric-comp-rot}
  \mathrm{MAE}^{\mathrm{comp}}_{\mathrm{rot}}
  = \Bigl\langle\; \bigl|\operatorname{wrap}\bigl(\widehat{\Delta\varphi}_{ik} - \widehat{\Delta\varphi}_{ij} - \widehat{\Delta\varphi}_{jk}\bigr)\bigr| \;\Bigr\rangle,
  \qquad \widehat{\Delta\varphi} = \operatorname{atan2}(\hat s, \hat c).
\end{equation}
Note \cref{eq:metric-comp-rot} composes the \emph{decoded angles}, not the
$(\cos,\sin)$ vectors, so it is insensitive to the head's output norm.

\emph{Angular error} itself, the headline rotation metric:
\begin{equation}\label{eq:metric-ang-mae}
  \mathrm{MAE}_\varphi
  = \Bigl\langle\, \bigl|\operatorname{wrap}\bigl(\operatorname{atan2}(\hat s_{ij}, \hat c_{ij}) - \Delta\varphi_{ij}\bigr)\bigr| \,\Bigr\rangle_{i \neq j},
\end{equation}
read against \cref{eq:ang-floor}.

\emph{Quad accuracy}, for the discrete arms only --- the fraction of pairs whose decoded
angle rounds to the correct quarter-turn:
\begin{equation}\label{eq:metric-quad-acc}
  \mathrm{Acc}_{\mathrm{quad}}
  = \Bigl\langle\, \mathbb{1}\bigl[\, |\operatorname{wrap}(\widehat{\Delta\varphi} - \Delta\varphi)| < \tfrac{\pi}{4} \,\bigr] \Bigr\rangle,
  \qquad \text{chance} = 25\% .
\end{equation}

\paragraph{The composite positive criterion.} Worth stating as a single displayed
condition, because it is the standard the whole of \cref{ch:pretraining} is judged
against:
\begin{equation}\label{eq:positive-criterion}
  \underbrace{\mathrm{MAE}_\varphi \ll \tfrac{2}{3}a}_{\text{beats the floor}}
  \;\;\wedge\;\;
  \underbrace{\mathrm{Id}_{\mathrm{rot}},\, \rho^{\mathrm{sym}}_{\mathrm{rot}},\, \mathrm{MAE}^{\mathrm{comp}}_{\mathrm{rot}} \to 0}_{\text{invariants hold}}
  \;\;\wedge\;\;
  \underbrace{\text{survives control (c)}}_{\text{not the signature}} .
\end{equation}


\ifSubfilesClassLoaded{\pagebreak\printbibliography}{}
\end{document}
